\documentclass[12pt,a4paper]{article}

\usepackage[utf8]{inputenc}
\usepackage[T1]{fontenc}
\usepackage[brazil]{babel}
\usepackage[a4paper,margin=2.cm]{geometry}
\usepackage{enumitem}
\setlist[enumerate]{itemsep=0.45em,topsep=0.35em}

\begin{document}

\begin{center}
  {\Large\bfseries PROVA SIMULADA -- NPOR -- VERSÃO 2}\par
  \vspace{0.4cm}
  Nome: \rule{9cm}{0.4pt}\par\vspace{0.3cm}
  Data: \rule{4cm}{0.4pt}\hfill Nota: \rule{3cm}{0.4pt}
\end{center}

\noindent\textbf{Duração sugerida: 90 minutos.} A prova oficial pode ter outro tempo e outros conteúdos; confirme as instruções da sua Organização Militar. Este simulado tem 13 questões objetivas e uma redação.

\medskip
\noindent\textbf{Instruções:} marque uma alternativa por questão. Na redação, escreva um texto dissertativo-argumentativo, com letra legível, respeitando o tema e o espaço disponível.

\subsection*{Português}
Leia o texto para responder às questões 1 e 2.

\begin{quote}
Uma praça bem cuidada pode aproximar os moradores e tornar o bairro mais agradável. Para isso, não basta construir equipamentos: é preciso conservar o espaço e ouvir quem o utiliza. Quando a comunidade participa, as soluções tendem a atender melhor às necessidades locais.
\end{quote}

\begin{enumerate}
  \item De acordo com o texto, o que ajuda a tornar as soluções para a praça mais adequadas?
  \begin{enumerate}[label=\alph*)]
    \item Impedir que os moradores utilizem o espaço.
    \item Priorizar somente a construção de equipamentos.
    \item Ouvir a comunidade e conservar o local.
    \item Transferir toda a responsabilidade aos visitantes.
  \end{enumerate}

  \item No trecho ``Para isso'', a expressão retoma a ideia de:
  \begin{enumerate}[label=\alph*)]
    \item aproximar moradores e tornar o bairro mais agradável.
    \item construir casas ao redor da praça.
    \item reduzir o número de espaços públicos.
    \item substituir a participação da comunidade.
  \end{enumerate}

  \item Assinale a frase com concordância verbal adequada.
  \begin{enumerate}[label=\alph*)]
    \item Havia poucos ônibus disponíveis naquele horário.
    \item Fazem três semanas que a obra começou.
    \item Os moradores participou da reunião.
    \item Existe muitas opções de transporte.
  \end{enumerate}

  \item Assinale a alternativa em que o uso da crase está correto.
  \begin{enumerate}[label=\alph*)]
    \item Os moradores chegaram à praça antes da reunião.
    \item A equipe começou à organizar os materiais.
    \item O ônibus chegou à tempo.
    \item O aviso foi entregue à todos os moradores.
  \end{enumerate}
\end{enumerate}

\subsection*{Matemática e raciocínio lógico}
\begin{enumerate}[resume]
  \item Uma turma arrecadou R\$ 80,00. Se gastou $\frac{3}{4}$ desse valor, quanto gastou?
  \begin{enumerate}[label=\alph*)]
    \item R\$ 20,00
    \item R\$ 40,00
    \item R\$ 60,00
    \item R\$ 75,00
  \end{enumerate}

  \item Se 6 trabalhadores concluem um serviço em 5 dias, em quantos dias 10 trabalhadores o concluiriam, mantendo o mesmo ritmo?
  \begin{enumerate}[label=\alph*)]
    \item 2 dias
    \item 3 dias
    \item 5 dias
    \item 8 dias
  \end{enumerate}

  \item No sistema $x+y=14$ e $x-y=4$, qual é o valor de $x$?
  \begin{enumerate}[label=\alph*)]
    \item 5
    \item 7
    \item 9
    \item 14
  \end{enumerate}

  \item Um terreno retangular mede 9 m de comprimento e 4 m de largura. Qual é seu perímetro?
  \begin{enumerate}[label=\alph*)]
    \item 13 m
    \item 26 m
    \item 36 m
    \item 40 m
  \end{enumerate}

  \item Um cubo tem arestas de 3 cm. Qual é seu volume?
  \begin{enumerate}[label=\alph*)]
    \item 9 cm$^3$
    \item 18 cm$^3$
    \item 27 cm$^3$
    \item 36 cm$^3$
  \end{enumerate}
\end{enumerate}

\subsection*{Conhecimentos gerais e institucionais}
\begin{enumerate}[resume]
  \item Em que ano foi proclamada a República no Brasil?
  \begin{enumerate}[label=\alph*)]
    \item 1822
    \item 1888
    \item 1889
    \item 1930
  \end{enumerate}

  \item Em qual país a Força Expedicionária Brasileira atuou durante a Segunda Guerra Mundial?
  \begin{enumerate}[label=\alph*)]
    \item Itália.
    \item Canadá.
    \item Japão.
    \item Portugal.
  \end{enumerate}

  \item Qual é o maior estado brasileiro em extensão territorial?
  \begin{enumerate}[label=\alph*)]
    \item Bahia.
    \item Amazonas.
    \item Pará.
    \item Mato Grosso.
  \end{enumerate}

  \item Qual posto vem imediatamente depois de Capitão na hierarquia do oficialato?
  \begin{enumerate}[label=\alph*)]
    \item Major.
    \item Primeiro-Tenente.
    \item Tenente-Coronel.
    \item Coronel.
  \end{enumerate}
\end{enumerate}

\newpage
\section*{Redação}
\noindent\textbf{Tema:} A cooperação e o senso de responsabilidade na construção de uma comunidade melhor.

\medskip
\noindent Escreva um texto dissertativo-argumentativo, apresentando seu ponto de vista, argumentos que o sustentem e uma conclusão. Sugestão: de 15 a 25 linhas. Não é necessário escrever um título.

\vspace{0.4cm}
\noindent\rule{\textwidth}{0.3pt}\par\vspace{0.45cm}
\noindent\rule{\textwidth}{0.3pt}\par\vspace{0.5cm}
\noindent\rule{\textwidth}{0.3pt}\par\vspace{0.5cm}
\noindent\rule{\textwidth}{0.3pt}\par\vspace{0.5cm}
\noindent\rule{\textwidth}{0.3pt}\par\vspace{0.5cm}
\noindent\rule{\textwidth}{0.3pt}\par\vspace{0.5cm}
\noindent\rule{\textwidth}{0.3pt}\par\vspace{0.5cm}
\noindent\rule{\textwidth}{0.3pt}\par\vspace{0.5cm}
\noindent\rule{\textwidth}{0.3pt}\par\vspace{0.5cm}
\noindent\rule{\textwidth}{0.3pt}\par\vspace{0.5cm}
\noindent\rule{\textwidth}{0.3pt}\par\vspace{0.5cm}
\noindent\rule{\textwidth}{0.3pt}\par\vspace{0.5cm}
\noindent\rule{\textwidth}{0.3pt}\par\vspace{0.5cm}
\noindent\rule{\textwidth}{0.3pt}\par\vspace{0.5cm}
\noindent\rule{\textwidth}{0.3pt}\par\vspace{0.5cm}
\noindent\rule{\textwidth}{0.3pt}\par\vspace{0.5cm}
\noindent\rule{\textwidth}{0.3pt}\par\vspace{0.5cm}
\noindent\rule{\textwidth}{0.3pt}\par\vspace{0.5cm}
\noindent\rule{\textwidth}{0.3pt}\par\vspace{0.5cm}
\noindent\rule{\textwidth}{0.3pt}\par\vspace{0.5cm}
\noindent\rule{\textwidth}{0.3pt}\par\vspace{0.5cm}
\noindent\rule{\textwidth}{0.3pt}\par

\end{document}